\section{Sampling and relative-entropy estimates}

\subsection{Discrepancy bounds}

We first record the elementary estimate which passes from star discrepancy to
histogram error.  The only point to note is the number of interval atoms in the
rotation partition.

\begin{lemma}[Interval partitions from star discrepancy]\label{lem:interval-partition-star-discrepancy}
Let \(\nu_N:=N^{-1}\sum_{t=0}^{N-1}\delta_{x_0+t\alpha}\) and let
\(\lambda\) be Lebesgue measure on \(\TT\).  If
\(\mathcal P=\{I_1,\dots,I_r\}\) is a finite partition of \(\TT\) into
half-open intervals, then
\[
\TV\bigl((\nu_N(I))_{I\in\mathcal P},(\lambda(I))_{I\in\mathcal P}\bigr)
\le r\,D_N^*(\alpha;x_0).
\]
\end{lemma}

\begin{proof}
For every half-open interval \(I\subset\TT\), possibly wrapping around the
origin,
\[
|\nu_N(I)-\lambda(I)|\le 2D_N^*(\alpha;x_0),
\]
because \(I\) is the difference, modulo complements if necessary, of two
anchored intervals \([0,a)\).  Therefore
\[
\TV\bigl((\nu_N(I))_{I\in\mathcal P},(\lambda(I))_{I\in\mathcal P}\bigr)
=\frac12\sum_{I\in\mathcal P}|\nu_N(I)-\lambda(I)|
\le rD_N^*(\alpha;x_0).
\]
\end{proof}

\begin{theorem}[Discrepancy control of microstate and folded laws]\label{thm:tv-certificate-hist}
For every irrational \(\alpha\), every \(\beta\in(0,1)\), every
\(x_0\in\TT\), and every \(m,N\ge1\),
\[
\TV(\hat\mu_{m,N},\mu_m)
\le 2m\,D_N^*(\alpha;x_0),
\qquad
\TV(\hat\pi_{m,N},\pi_m^{\mathrm{fold}})
\le 2m\,D_N^*(\alpha;x_0).
\]
In the Sturmian slice \(\beta=\alpha\), these improve to
\[
\TV(\hat\mu_{m,N},\mu_m)
\le (m+1)D_N^*(\alpha;x_0),
\qquad
\TV(\hat\pi_{m,N},\pi_m^{\mathrm{fold}})
\le (m+1)D_N^*(\alpha;x_0).
\]
\end{theorem}

\begin{proof}
The length-\(m\) coding map
\[
x\longmapsto \omega_m(x)
=\bigl(\mathbf 1_{[0,\beta)}(x+j\alpha)\bigr)_{j=0}^{m-1}
\]
can change only when \(x\) crosses one of the boundary points
\[
\Gamma_m(\alpha,\beta)
:=\{-j\alpha,\ \beta-j\alpha \pmod 1:\ 0\le j<m\}.
\]
After deleting repetitions and arranging these points cyclically, they cut
\(\TT\) into a half-open interval partition \(\mathcal P_m\) with
\(|\mathcal P_m|\le 2m\).  The convention \([0,\beta)\) makes
\(\omega_m\) constant on each atom of this partition, so the microstate laws
\(\hat\mu_{m,N}\) and \(\mu_m\) are pushforwards of the two partition laws
\((\nu_N(I))_{I\in\mathcal P_m}\) and \((\lambda(I))_{I\in\mathcal P_m}\).
Total variation cannot increase under a deterministic pushforward.  Hence
Lemma~\ref{lem:interval-partition-star-discrepancy} gives
\[
\TV(\hat\mu_{m,N},\mu_m)
\le |\mathcal P_m|D_N^*(\alpha;x_0)
\le 2mD_N^*(\alpha;x_0).
\]
Since \(\hat\pi_{m,N}=(\Fold_m)_*\hat\mu_{m,N}\) and
\(\pi_m^{\mathrm{fold}}=(\Fold_m)_*\mu_m\), another application of
monotonicity under deterministic pushforward gives the folded estimate.

If \(\beta=\alpha\), then
\[
\Gamma_m(\alpha,\alpha)
=\{-j\alpha,\ \alpha-j\alpha \pmod 1:\ 0\le j<m\}
=\{-(j-1)\alpha\pmod 1:\ 0\le j\le m\}.
\]
These \(m+1\) points are distinct because \(\alpha\) is irrational.  Thus the
same argument applies with \(|\mathcal P_m|=m+1\), proving the sharpened
Sturmian bounds.
\end{proof}

\begin{remark}
Thus any arithmetic or numerical bound on \(D_N^*(\alpha;x_0)\) gives the same
order of control for the folded histogram.  No independence approximation is
being used.
\end{remark}

\subsection{Cross-resolution residuals}

Let $r_{m+1\to m}:\Omega_{m+1}\to\Omega_m$ denote prefix truncation,
$r_{m+1\to m}(u_1\cdots u_{m+1})=u_1\cdots u_m$, and let
$\rho_{m+1\to m}:X_{m+1}\to X_m$ be the analogous truncation on folded words.
The following residual measures the failure of folding and truncation to
commute between adjacent window lengths:
Define the projective residual
\[
E_{m,N}
:=
\TV\!\left(\hat\pi_{m,N},(\rho_{m+1\to m})_*\hat\pi_{m+1,N}\right).
\]

\begin{theorem}[Boundary criterion for the canonical fold]\label{thm:multiscale-residual-by-boundary}
For \(u=u_1\cdots u_{m+1}\in\Omega_{m+1}\), let
\(a:=u_1\cdots u_m\). Then
\[
\rho_{m+1\to m}\Fold_{m+1}(u)\ne \Fold_m(a)
\]
if and only if
\[
u_{m+1}=1
\quad\text{and}\quad
N_m(a)\ge F_{m+1}.
\]
Consequently, with
\[
B_m:=\{u\in\Omega_{m+1}:u_{m+1}=1,\ N_m(u_1\cdots u_m)\ge F_{m+1}\},
\]
one has
\[
E_{m,N}\le \hat\mu_{m+1,N}(B_m).
\]
\end{theorem}

\begin{proof}
Write \(n:=N_m(a)\). If \(u_{m+1}=0\), then
\(N_{m+1}(u)=n\), so the first \(m\) Zeckendorf digits are the same at levels
\(m\) and \(m+1\).

Now suppose \(u_{m+1}=1\). If \(n<F_{m+1}\), the Zeckendorf expansion of \(n\)
uses no digit in positions \(m\) or \(m+1\). Adding \(F_{m+2}\) therefore
places an isolated \(1\) in position \(m+1\) and leaves the first \(m\) digits
unchanged.

It remains to show that the first \(m\) digits cannot remain unchanged when
\(n\ge F_{m+1}\). Write \(p:=N_m(\Fold_m(a))\). By
Proposition~\ref{prop:canonical-fold-basic}, \(n=p+\varepsilon F_{m+2}\) with
\(\varepsilon\in\{0,1\}\). Also \(N_{m+1}(u)<F_{m+4}\), so no Zeckendorf digit
beyond position \(m+2\) can occur. If \(F_{m+1}\le n<F_{m+2}\), then
\(\varepsilon=0\), \(p=n\), and the Zeckendorf expansion of \(p\) has digit
\(1\) in position \(m\). Keeping the first \(m\) digits after adding
\(F_{m+2}\) would force the digit in position \(m+1\) also to be \(1\), giving
an illegal adjacent pair. No alternative choice of higher digits has the same
value, because the only representation of
\[
p+F_{m+2}=p+bF_{m+2}+cF_{m+3},
\qquad b,c\in\{0,1\},
\]
with the first \(m\) digits fixed is \(b=1,c=0\).

If \(n\ge F_{m+2}\), then \(\varepsilon=1\). Were the first \(m\) digits
unchanged after adding \(F_{m+2}\), the higher digits would have to satisfy
\[
2F_{m+2}=bF_{m+2}+cF_{m+3},
\qquad b,c\in\{0,1\},
\]
which is impossible. Hence the first \(m\) digits change exactly in the stated
case.

For deterministic maps \(A(u)=\Fold_m(r_{m+1\to m}u)\) and
\(B(u)=\rho_{m+1\to m}\Fold_{m+1}(u)\), the total variation distance between
the empirical laws of \(A(u)\) and \(B(u)\) is bounded by the empirical mass of
\(\{A\ne B\}\). The criterion identifies this set with \(B_m\).
\end{proof}

\begin{proposition}[Residuals relative to the Parry measure]\label{prop:multiscale-residual-parry-bound}
Let \(\pi_m\) be the length-\(m\) Parry block law on \(X_m\). Then, for every
\(m,N\ge1\),
\[
E_{m,N}
\le
\TV(\hat\pi_{m,N},\pi_m)+\TV(\hat\pi_{m+1,N},\pi_{m+1}).
\]
\end{proposition}

\begin{proof}
The Parry family is projectively consistent:
\[
\pi_m=(\rho_{m+1\to m})_*\pi_{m+1}.
\]
Hence
\[
\begin{aligned}
E_{m,N}
&=
\TV\!\left(\hat\pi_{m,N},(\rho_{m+1\to m})_*\hat\pi_{m+1,N}\right)\\
&\le
\TV(\hat\pi_{m,N},\pi_m)
+
\TV\!\left(\pi_m,(\rho_{m+1\to m})_*\hat\pi_{m+1,N}\right)\\
&=
\TV(\hat\pi_{m,N},\pi_m)
+
\TV\!\left((\rho_{m+1\to m})_*\pi_{m+1},
(\rho_{m+1\to m})_*\hat\pi_{m+1,N}\right).
\end{aligned}
\]
Total variation does not increase under pushforward by a deterministic map, so
\[
\TV\!\left((\rho_{m+1\to m})_*\pi_{m+1},
(\rho_{m+1\to m})_*\hat\pi_{m+1,N}\right)
\le
\TV(\pi_{m+1},\hat\pi_{m+1,N}).
\]
Substituting this estimate proves the claim.
\end{proof}

\begin{corollary}[Residual comparison bound]
\label{cor:multiscale-residual-comparison-bound}
For every irrational \(\alpha\), every \(\beta\in(0,1)\), and every
\(m,N\ge1\),
\[
E_{m,N}
\le
(4m+2)D_N^*(\alpha;x_0)
+
\TV(\pi_m^{\mathrm{fold}},\pi_m)
+
\TV(\pi_{m+1}^{\mathrm{fold}},\pi_{m+1}).
\]
In the Sturmian slice \(\beta=\alpha\), one may replace \(4m+2\) by
\(2m+3\).
\end{corollary}

\begin{proof}
By Proposition~\ref{prop:multiscale-residual-parry-bound},
\[
E_{m,N}
\le
\TV(\hat\pi_{m,N},\pi_m)+\TV(\hat\pi_{m+1,N},\pi_{m+1}).
\]
Apply the triangle inequality at levels \(m\) and \(m+1\):
\[
\TV(\hat\pi_{\ell,N},\pi_\ell)
\le
\TV(\hat\pi_{\ell,N},\pi_\ell^{\mathrm{fold}})
+
\TV(\pi_\ell^{\mathrm{fold}},\pi_\ell),
\qquad \ell\in\{m,m+1\}.
\]
Now use Theorem~\ref{thm:tv-certificate-hist}. In the general window this gives
\[
\TV(\hat\pi_{m,N},\pi_m^{\mathrm{fold}})
+
\TV(\hat\pi_{m+1,N},\pi_{m+1}^{\mathrm{fold}})
\le
(2m+2m+2)D_N^*(\alpha;x_0),
\]
while in the Sturmian slice it gives
\[
\TV(\hat\pi_{m,N},\pi_m^{\mathrm{fold}})
+
\TV(\hat\pi_{m+1,N},\pi_{m+1}^{\mathrm{fold}})
\le
((m+1)+(m+2))D_N^*(\alpha;x_0).
\]
Substituting these two estimates proves the corollary.
\end{proof}
\subsection{Relative-entropy bounds}

\subsubsection{From total variation to relative entropy}

To pass from total variation to relative entropy one needs a lower bound on the
smallest atom of the target distribution.

\begin{lemma}[Relative entropy from total variation under strict positivity]\label{lem:kl-from-tv-qmin}
Let \(p\) and \(q\) be probability laws on a finite set \(E\), with
\(\supp(p)\subset\supp(q)\), and put
\[
q_{\min}:=\min_{x\in\supp(q)}q(x)>0.
\]
Then
\[
\KL(p\|q)\le \frac{4}{q_{\min}}\TV(p,q)^2.
\]
\end{lemma}

\begin{proof}
The standard inequality \(\KL(p\|q)\le \chi^2(p,q)\) gives
\[
\KL(p\|q)
\le \sum_{x\in\supp(q)}\frac{(p(x)-q(x))^2}{q(x)}
\le \frac{1}{q_{\min}}\sum_{x\in E}|p(x)-q(x)|^2.
\]
Since \(\sum_x |p(x)-q(x)|=2\TV(p,q)\), the last expression is at most
\(4q_{\min}^{-1}\TV(p,q)^2\).  This is the finite-alphabet
Pinsker--Csisz\'ar--type comparison used below
\cite{Csiszar1975IdivergenceGeometry,SasonVerdu2016fDivergence}.
\end{proof}

\subsubsection{Lower bounds on the smallest rotation atom}

\begin{definition}[Diophantine lower bound]\label{def:diophantine-lower-bound}
For an irrational \(\alpha\), set
\[
c(\alpha):=\inf_{k\ge1} k\|k\alpha\|,
\]
where \(\|x\|\) denotes distance to the nearest integer. Thus
\(\alpha\) is badly approximable if and only if \(c(\alpha)>0\).
This is the one-dimensional instance of Schmidt's badly-approximable
constant convention \cite{Schmidt1969BadlyApproximableSystems}.
\end{definition}

\begin{definition}[General window partition scale]\label{def:general-window-partition-scale}
For \(m\ge1\), let
\[
\Delta_m(\alpha,\beta)
:=\{-j\alpha,\ \beta-j\alpha \pmod 1:\ 0\le j<m\}\subset\TT .
\]
Let \(\mathcal P_m(\alpha,\beta)\) be the half-open circle partition whose
endpoints are the distinct points of \(\Delta_m(\alpha,\beta)\), taken in
cyclic order and oriented counterclockwise.  Define
\[
\lambda_m(\alpha,\beta)
:=
\min_{I\in\mathcal P_m(\alpha,\beta)}\leb(I).
\]
\end{definition}

\begin{proposition}[General window atom separation]\label{prop:general-window-atom-separation}
For every irrational \(\alpha\), every \(\beta\in(0,1)\), and every
\(m\ge1\), each positive atom of \(\mu_m\) and of
\(\pi_m^{\mathrm{fold}}\) has mass at least
\(\lambda_m(\alpha,\beta)\).  Moreover
\[
\lambda_m(\alpha,\beta)
=
\min\Bigl(
\{\|k\alpha\|:1\le k\le m-1\}
\cup
\{\|\beta+r\alpha\|:\ |r|\le m-1,\ \beta+r\alpha\notin\mathbb Z\}
\Bigr),
\]
with the first set omitted when \(m=1\).  Consequently, if
\(c(\alpha)>0\) and
\[
s_m(\alpha,\beta)
:=
\min_{\substack{|r|\le m-1\\ \beta+r\alpha\notin\mathbb Z}}
\|\beta+r\alpha\|,
\]
then
\[
\lambda_m(\alpha,\beta)\ge
\min\left\{\frac{c(\alpha)}{m},\,s_m(\alpha,\beta)\right\}.
\]
If, in addition,
\[
c_\beta(\alpha):=
\inf_{\substack{r\in\mathbb Z\\ \beta+r\alpha\notin\mathbb Z}}
(|r|+1)\|\beta+r\alpha\|>0,
\]
then
\[
\lambda_m(\alpha,\beta)
\ge
\frac{\min\{c(\alpha),c_\beta(\alpha)\}}{m}.
\]
\end{proposition}

\begin{proof}
The discontinuity set of the word map \(x\mapsto\omega_m(x)\) is contained in
\(\Delta_m(\alpha,\beta)\).  On each half-open atom of
\(\mathcal P_m(\alpha,\beta)\) the word is constant.  The half-open convention
is compatible with the coding interval \([0,\beta)\): at a point
\(-j\alpha\) the \(j\)-th coordinate takes the right-hand value \(1\), while at
\(\beta-j\alpha\) it takes the right-hand value \(0\).  Thus every observed
boundary word also occurs on the adjacent half-open atom of positive length.
Therefore each cylinder \(C_m(a)\) is a union of atoms of
\(\mathcal P_m(\alpha,\beta)\), and every positive cylinder mass
\(\mu_m(a)\) is at least \(\lambda_m(\alpha,\beta)\).  A positive folded atom
is a sum of positive cylinder masses over one fiber of \(\Fold_m\), so it is
also at least \(\lambda_m(\alpha,\beta)\).

The minimum atom length of a finite circle partition is the minimum positive
circular distance between two distinct endpoints.  Two endpoints from the same
family \(\{-j\alpha\}\) or \(\{\beta-j\alpha\}\) differ by \(k\alpha\), with
\(1\le |k|\le m-1\).  Endpoints from different families differ by
\(\beta+r\alpha\), with \(|r|\le m-1\).  Removing the zero cross-differences
coming from coincident endpoints gives the displayed formula for
\(\lambda_m(\alpha,\beta)\).

For \(1\le k\le m-1\),
\[
\|k\alpha\|\ge \frac{c(\alpha)}{k}\ge \frac{c(\alpha)}{m}.
\]
The cross-family distances are bounded below by \(s_m(\alpha,\beta)\), which
proves the finite-scale bound.  If \(c_\beta(\alpha)>0\), then for
\(|r|\le m-1\),
\[
\|\beta+r\alpha\|
\ge \frac{c_\beta(\alpha)}{|r|+1}
\ge \frac{c_\beta(\alpha)}{m}.
\]
Combining this with the same-family estimate gives the final inequality.
\end{proof}

\begin{proposition}[Finite-scale lower atom bound]\label{prop:qmin-lowerbound-markov}
In the Sturmian slice \(\beta=\alpha\), write \(q_{\min}(m)\) for the least
positive atom of \(\mu_m\). If \(\alpha\) has Diophantine constant
\(c(\alpha)\), then
\[
q_{\min}(m)\ge \frac{c(\alpha)}{m}.
\]
For \(\alpha=\varphi^{-1}\), the all-\(k\) constant \(c(\alpha)=\varphi^{-2}\)
is valid, and therefore
\[
q_{\min}(m)\ge \frac{\varphi^{-2}}{m}.
\]
\end{proposition}

\begin{proof}
In the Sturmian slice, the positive atoms of \(\mu_m\) are the gaps between the
points \(\{(1-j)\alpha:0\le j\le m\}\) on the circle. Every such gap has length
at least
\[
\min_{1\le k\le m}\|k\alpha\|
\ge \frac{c(\alpha)}{m},
\]
which proves the first claim.

For \(\alpha=\varphi^{-1}=[0;1,1,\ldots]\), the continued-fraction
denominators in the convention \(q_{-1}=0,q_0=1\) are
\[
q_j=F_{j+1}\qquad(j\ge0),
\]
where \(F_0=0,F_1=1\).  Since
\(\varphi^{-1}=1/\varphi\) and \(F_n=(\varphi^n-(-\varphi)^{-n})/\sqrt5\),
one has, for \(n\ge2\),
\[
F_n\varphi^{-1}-F_{n-1}=(-1)^{n+1}\varphi^{-n},
\qquad\text{hence}\qquad
\|F_n\varphi^{-1}\|=\varphi^{-n}.
\]
Equivalently, after substituting \(n=j+1\), this is Cassels' denominator
estimate
\[
\|q_j\alpha\|=\|F_{j+1}\alpha\|=\varphi^{-j-1}
\]
in this Fibonacci indexing for \(j\ge1\) \cite[Ch.~I]{Cassels1957}.  The
initial denominator \(q_0=F_1=1\) instead gives
\(\|q_0\alpha\|=\alpha=\varphi^{-1}\), and hence the same lower bound for
\(k\|k\alpha\|\).  For \(k\ge2\), choose \(n\ge2\) with
\(F_n\le k<F_{n+1}\).  Since \(q_{n-1}=F_n\) and \(q_n=F_{n+1}\), the
best-approximation property of continued-fraction denominators says that for
every integer \(1\le k<q_n=F_{n+1}\) with \(k\ne q_{n-1}\),
\[
\|k\alpha\|\ge \|q_{n-1}\alpha\|=\|F_n\alpha\|=\varphi^{-n}.
\]
The same lower bound is an equality when \(k=F_n\). Hence
\[
k\|k\alpha\|\ge F_n\varphi^{-n}\ge \varphi^{-2},
\]
where the last inequality uses \(F_n\ge\varphi^{n-2}\). This proves the stated
uniform choice of \(c(\alpha)\).
\end{proof}

\subsubsection{Microstate and folded relative-entropy bounds}

\begin{theorem}[General-window relative-entropy certificate]\label{thm:general-window-relative-entropy-certificate}
For every irrational \(\alpha\), every \(\beta\in(0,1)\), and every
\(m,N\ge1\),
\[
\KL(\hat\mu_{m,N}\|\mu_m)
\le
\frac{16m^2}{\lambda_m(\alpha,\beta)}
D_N^*(\alpha;x_0)^2
\]
and
\[
\KL(\hat\pi_{m,N}\|\pi_m^{\mathrm{fold}})
\le
\frac{16m^2}{\lambda_m(\alpha,\beta)}
D_N^*(\alpha;x_0)^2 .
\]
Consequently,
\[
\KL(\hat\mu_{m,N}\|\mu_m),\
\KL(\hat\pi_{m,N}\|\pi_m^{\mathrm{fold}})
\le
\frac{16m^2}
{\min\{c(\alpha)/m,s_m(\alpha,\beta)\}}
D_N^*(\alpha;x_0)^2
\]
whenever \(c(\alpha)>0\) and the finite-scale window separation in
Proposition~\ref{prop:general-window-atom-separation} is available.  If
\(c(\alpha)>0\) and \(c_\beta(\alpha)>0\), then both relative entropies are at
most
\[
\frac{16m^3}{\min\{c(\alpha),c_\beta(\alpha)\}}
D_N^*(\alpha;x_0)^2 .
\]
\end{theorem}

\begin{proof}
By Proposition~\ref{prop:general-window-atom-separation}, the least positive
atom of \(\mu_m\) and the least positive atom of
\(\pi_m^{\mathrm{fold}}\) are both at least
\(\lambda_m(\alpha,\beta)\).  The same proposition also shows that the
half-open coding convention puts every empirical word, including a boundary
word, inside a cylinder of positive Lebesgue mass.  Hence
\(\supp(\hat\mu_{m,N})\subset\supp(\mu_m)\), and after applying the
deterministic fold,
\(\supp(\hat\pi_{m,N})\subset\supp(\pi_m^{\mathrm{fold}})\).

Apply Lemma~\ref{lem:kl-from-tv-qmin} to the microstate law and to the folded
law, and then use the general-window part of
Theorem~\ref{thm:tv-certificate-hist}:
\[
\TV(\hat\mu_{m,N},\mu_m),\
\TV(\hat\pi_{m,N},\pi_m^{\mathrm{fold}})
\le 2mD_N^*(\alpha;x_0).
\]
This gives the two displayed \(\lambda_m(\alpha,\beta)\)-bounds.  The
remaining estimates are obtained by substituting the separation lower bounds
from Proposition~\ref{prop:general-window-atom-separation}.
\end{proof}

The next proposition and corollaries keep the sharper \((m+1)\)-atom Sturmian
constant and the explicit bad-approximation lower bound used later in the
Parry comparison.

\begin{proposition}[Sturmian microstate relative-entropy certificate]\label{prop:sturmian-microstate-kl-certificate}
Assume \(\beta=\alpha\). Then
\[
\KL(\hat\mu_{m,N}\|\mu_m)
\le
\frac{4(m+1)^2}{q_{\min}(m)}D_N^*(\alpha;x_0)^2.
\]
If \(\alpha\) is badly approximable, with Diophantine constant
\(c(\alpha)>0\), then
\[
\KL(\hat\mu_{m,N}\|\mu_m)
\le
\frac{4m(m+1)^2}{c(\alpha)}D_N^*(\alpha;x_0)^2.
\]
\end{proposition}

\begin{proof}
Every atom counted by \(\hat\mu_{m,N}\) is a length-\(m\) word realized by the
same Sturmian coding, so
\(\supp(\hat\mu_{m,N})\subset S_m(\alpha)=\supp(\mu_m)\).  By the sharpened
Sturmian part of Theorem~\ref{thm:tv-certificate-hist},
\[
\TV(\hat\mu_{m,N},\mu_m)\le (m+1)D_N^*(\alpha;x_0).
\]
Applying Lemma~\ref{lem:kl-from-tv-qmin} with
\(q_{\min}=q_{\min}(m)\) therefore gives
\[
\KL(\hat\mu_{m,N}\|\mu_m)
\le
\frac{4}{q_{\min}(m)}(m+1)^2D_N^*(\alpha;x_0)^2.
\]
If \(c(\alpha)>0\), then Proposition~\ref{prop:qmin-lowerbound-markov} yields
\(q_{\min}(m)\ge c(\alpha)/m\), and substituting this lower bound proves the
second estimate.
\end{proof}

\begin{corollary}[Golden-slope microstate improvement]\label{cor:rotation-microstate-kl-certificate}
Assume \(\beta=\alpha=\varphi^{-1}\).  Then
\[
\KL(\hat\mu_{m,N}\|\mu_m)
\le
4\varphi^2\,m(m+1)^2 D_N^*(\alpha;x_0)^2.
\]
\end{corollary}

\begin{proof}
Apply Proposition~\ref{prop:sturmian-microstate-kl-certificate} and then use
the golden lower bound
\(q_{\min}(m)\ge \varphi^{-2}/m\) from
Proposition~\ref{prop:qmin-lowerbound-markov}.
\end{proof}

\begin{lemma}[KL monotonicity under pushforward]\label{lem:kl-monotone-pushforward}
Let $F:S\to T$ be a map between finite sets and let $p,q$ be probability
measures on $S$ with $p\ll q$.  Then
\[
\KL(F_*p\|F_*q)\le \KL(p\|q).
\]
\end{lemma}

\begin{proof}
This is the data-processing inequality for relative entropy under deterministic
coarse graining \cite{CoverThomas2006}.
\end{proof}

\begin{corollary}[Folded relative-entropy bound]\label{cor:rotation-folded-kl-certificate}
Assume \(\beta=\alpha\) and let \(\alpha\) be badly approximable. Then
\[
\KL(\hat\pi_{m,N}\|\pi_m^{\mathrm{fold}})
\le
\frac{4(m+1)^2}{q_{\min}(m)}D_N^*(\alpha;x_0)^2.
\]
For the golden slope,
\[
\KL(\hat\pi_{m,N}\|\pi_m^{\mathrm{fold}})
\le
4\varphi^2\,m(m+1)^2 D_N^*(\alpha;x_0)^2.
\]
\end{corollary}

\begin{proof}
Since \(\hat\pi_{m,N}=(\Fold_m)_*\hat\mu_{m,N}\) and
\(\pi_m^{\mathrm{fold}}=(\Fold_m)_*\mu_m\), data processing
(Lemma~\ref{lem:kl-monotone-pushforward}) gives
\[
\KL(\hat\pi_{m,N}\|\pi_m^{\mathrm{fold}})
\le
\KL(\hat\mu_{m,N}\|\mu_m).
\]
The \(q_{\min}(m)\)-bound therefore follows from
Proposition~\ref{prop:sturmian-microstate-kl-certificate}, and the
golden-slope bound follows from
Corollary~\ref{cor:rotation-microstate-kl-certificate}.
\end{proof}
